\section*{Chapter \ref{ch:mx:build-an-agent-based-market} --- Build: an Agent-Based Market}

\begin{solution}{exo:mx:build-an-agent-based-market:1}
$(0.212-0.158)/0.046=1.17$ standard deviations, which adds $1.17^2=1.4$ to the distance (a quarter of the fresh distance of 5.4).
\end{solution}

\begin{solution}{exo:mx:build-an-agent-based-market:2}
As $\tau^{-(\alpha-1)}$: $\tau^{-0.5}$ for a tail of 1.5 and $\tau^{-1.5}$ for 2.5. Only the first is long memory ($\alpha<2$: the autocorrelations
are not summable); the calibrated tail of 2.5 gives short memory.
\end{solution}

\begin{solution}{exo:mx:build-an-agent-based-market:3}
Each fact's scaled gap has variance $1/6+1/8=0.29$, so the expected distance is $7\times0.29=2.0$. The floor, 6 fresh target hours against the
8-hour mean, is 1.7, close to the same expectation; the calibrated population's 5.4 is about 2.7 times that:
a misfit beyond sampling noise, mostly in clustering, depth and lag-1 sign memory.
\end{solution}

\begin{solution}{exo:mx:build-an-agent-based-market:4}
Two reasons. The mean of three hours carries more simulation noise than the mean of six ($7\times(1/3+1/8)=3.2$ expected against 2.0), and the point was
chosen because its three hours happened to land closest among 27, so its in-sample distance also reflects which hours they were: in the calibration
hours its kurtosis was 3.11, on fresh hours 1.90. The fresh distance is the honest one.
\end{solution}

\begin{solution}{exo:mx:build-an-agent-based-market:5}
Without them nothing pushes the price between the noise orders, which the providers' deep book absorbs: the mid sits still for most 30-second
intervals and moves only when a run of noise orders or a chartist burst empties a level. Returns that are mostly zero with rare jumps have a very large
kurtosis. The fundamentalists are what makes the price follow $V$ in small steps.
\end{solution}

\begin{solution}{exo:mx:build-an-agent-based-market:6}
Its fills are adversely selected: the fundamentalists trade against it when $V$ has moved, so without skew its inventory grows the wrong way and
stays there while the price keeps moving against it. Skew makes its quotes attract the trades that reduce the inventory, so it holds less (1.6 lots root mean square at a skew of 0.3, against 14.5) for less time and earns the spread on flow that turns over.
\end{solution}

\begin{solution}{exo:mx:build-an-agent-based-market:7}
Yes. On the five facts the best grid point becomes a noise rate of 1.5, a run tail of 2.5 and a chartists' rate of 0.1, at a distance of 5.1 on the
calibration hours; the chapter's point comes second at 7.7. The faster noise flow thins the book (depth about 9 lots against the target's 24) and widens
the spread (about 1.08 ticks against 1.04, more than two target standard deviations): the full calibration rejected it for its book, which the five
facts no longer see. The facts chosen decide the answer.
\end{solution}

\begin{solution}{exo:mx:build-an-agent-based-market:8}
Matching facts does not identify mechanisms: different populations can produce the same seven moments (here removing the chartists moved no fact
significantly, so any chartist share near the calibrated one fits as well), the facts left unmatched show the population lacks something the target
has, and the target is itself a model. The calibration says the population is good enough for some questions, not what real traders are.
\end{solution}

\begin{solution}{pb:mx:build-an-agent-based-market:1}
\textbf{1.} See the definitions.
\textbf{2.} Lux and Marchesi: interactions of fundamentalists and chartists produce fat tails and clustered volatility from news that has neither.
Farmer, Patelli and Zovko: random order placement explains 96\% of the cross-sectional variance of spreads and 76\% of that of price diffusion on 11
London stocks with one free parameter.
\textbf{3.} One-lot \omterm{def:mx:the-limit-order-book:orders}{limit orders} on both sides, half one tick from the opposite best and the rest up to ten ticks behind it, each cancelled at a
constant rate; \omterm{def:mx:the-limit-order-book:orders}{market orders} at a constant rate with signs in Pareto runs; \omterm{def:mx:the-limit-order-book:orders}{market orders} toward the hidden value at a rate proportional to the gap;
\omterm{def:mx:the-limit-order-book:orders}{market orders} along the 30-second trend at a rate proportional to its size.
\textbf{4.} They alone see $V$.
\textbf{5.} So that runs that differ by one agent share them (common random numbers).
\textbf{6.} See the definition; two lots at the others' best quotes (one tick inside when their spread exceeds two ticks), shifted by the skew times the
inventory, and no bid (offer) beyond 20 lots long (short).
\textbf{7.} Each new bid was set below its own previous bid when long, so the bid walked down refresh after refresh.
\textbf{8.} 14.5, 4.2 and 1.6 lots; $-116$, 448 and 545 USD an hour for skews 0, 0.1 and 0.3.
\textbf{9.} Power-law large orders split into constant children give sign autocorrelation decaying as $\tau^{-(\alpha-1)}$; the noise takers' sign runs
have Pareto lengths with tail $\alpha$.
\textbf{10.} An excess from the runs falling with slope $-1.15$ over lags 1 to 7 (the law says $-1.5$), gone by lag 10; the rest comes from the
fundamentalists.
\textbf{11.} See the definition; so that differences between grid points are not differences of luck.
\textbf{12.} Eight hours of \texttt{firm.tape} without news; kurtosis and lag-1 autocorrelation of 30-second returns, clustering, sign autocorrelation at
lags 1 and 10, spread, depth.
\textbf{13.} Noise 0.9, 1.2, 1.5; run tail none, 1.5, 2.5; chartists 0, 0.05, 0.1; best 1.2, 2.5, 0.05.
\textbf{14.} \emph{Named result}: after calibration the distance is 10.1 on the three calibration hours and 5.4 on six fresh hours, against 27.9 for
the starting point and a floor of 1.7 for the target market itself; the sign runs carry the lag-1 sign memory ($-2.77$, standard error 0.15), the fundamentalists the
price's movement (kurtosis $+8.6$ without them), the market maker the spread ($+9.4$), and the chartists nothing significant ($+1.11$ with a standard error of 0.55).
\textbf{15.} Volatility clustering, depth and lag-1 sign memory: the target has a random activity level and Hawkes-clustered flow the population lacks.
\textbf{16.} Interactive agent configurations were more realistic than replay, and more so with a fundamental from historical data.
\textbf{17.} A replayed tape does not react to the algorithm's orders, so impact and the others' response are missing.
\textbf{18.} Several populations match the same facts; an agent type can be removed with little change.
\textbf{19.} Use it comparatively, algorithm against algorithm with the same seeds, and check the conclusions do not hinge on the facts it misses.
\textbf{20.} To answer ``what if'' questions about mechanisms in a market that reacts.
\end{solution}

\begin{solution}{iq:mx:build-an-agent-based-market:1}
Agents that place and cancel limit and \omterm{def:mx:the-limit-order-book:orders}{market orders} at random rates, with no strategy; with the order-flow rates as inputs it predicts spreads and
price diffusion across stocks well (Farmer, Patelli and Zovko), showing how much the mechanism alone explains.
\end{solution}

\begin{solution}{iq:mx:build-an-agent-based-market:2}
Choose the facts that matter for the question (returns, order flow, book), measure them on the stock with their sampling variability, and minimise a
weighted distance over the population's parameters by simulated moments with common random numbers; check the fit on fresh simulated days and on
facts not used in the fit, and compare the distance with what sampling noise alone gives.
\end{solution}

\begin{solution}{iq:mx:build-an-agent-based-market:3}
One event queue ordered by time and a tie-breaking sequence, every random stream seeded per agent and per purpose, no wall clock or thread timing in
the logic; then the two algorithms face the same exogenous flow and their difference is not noise.
\end{solution}

\begin{solution}{iq:mx:build-an-agent-based-market:4}
To make the fills that reduce inventory more likely and those that add to it less, which cuts inventory risk and adverse selection; enough that the
inventory mean-reverts within the holding time the desk accepts, which in the chapter's market meant 0.1 to 0.3 ticks per lot.
\end{solution}

\begin{solution}{iq:mx:build-an-agent-based-market:5}
It learns the simulator's quirks: agents that do not adapt to it, liquidity that refills by fixed rules, missing facts such as clustering; it may
exploit patterns absent in reality. Randomise the population, test on held-out configurations and on real data before trading.
\end{solution}

\begin{solution}{iq:mx:build-an-agent-based-market:6}
Which facts the simulator matches and misses, whether the algorithm's advantage survives other plausible populations and seeds, whether its behaviour
in stressed scenarios (halts, gaps, one-sided flow) was tested, and a plan for a limited live trial with controls.
\end{solution}
